% ============================================================================
% Chapter 14: Computer Security Basics
% ============================================================================

\chapter{Computer Security Basics}
\label{ch:security}

\begin{objectives}
  \item Understand why computer security is important
  \item Create strong passwords and recognise weak ones
  \item Identify common threats: viruses, malware, phishing
  \item Use Windows security features to stay safe
\end{objectives}

\bytetip[tip]{Your computer is like your house --- you wouldn't leave the front door wide open, right? Computer security is about \textbf{locking the digital doors}!}

\section{Why Security Matters}
\label{sec:why-security}
\index{security}

Every day, millions of computers are targeted by hackers, viruses, and scams. Good security habits protect your personal information, school work, and the computer itself.

\section{Password Strength}
\label{sec:passwords}
\index{password!strength}

Not all passwords are equal. Here's how to judge password strength:

% --- Figure 14.1: Password Strength Meter ---
\begin{figure}[H]
\centering
\begin{tikzpicture}
  % Strength bar
  \fill[WarnRed] (0,0) rectangle (3,0.8);
  \fill[NoteOrange] (3,0) rectangle (6,0.8);
  \fill[FunYellow!80!black] (6,0) rectangle (9.5,0.8);
  \fill[LabGreen] (9.5,0) rectangle (13,0.8);
  
  % Labels
  \node[font=\small\bfseries\sffamily, text=white] at (1.5,0.4) {Weak};
  \node[font=\small\bfseries\sffamily, text=white] at (4.5,0.4) {Fair};
  \node[font=\small\bfseries\sffamily, text=white] at (7.75,0.4) {Good};
  \node[font=\small\bfseries\sffamily, text=white] at (11.25,0.4) {Strong};
  
  % Example passwords with pointers
  \draw[WarnRed, line width=1.5pt, ->] (1.0,-0.3) -- (1.0,-1.0);
  \node[below, font=\scriptsize\sffamily, text=WarnRed, align=center] at (1.0,-1.0) {\texttt{abc}\\Cracked in seconds};
  
  \draw[NoteOrange, line width=1.5pt, ->] (4.5,-0.3) -- (4.5,-1.0);
  \node[below, font=\scriptsize\sffamily, text=NoteOrange, align=center] at (4.5,-1.0) {\texttt{pass123}\\Cracked in minutes};
  
  \draw[LabGreen, line width=1.5pt, ->] (11.0,-0.3) -- (11.0,-1.0);
  \node[below, font=\scriptsize\sffamily, text=LabGreen, align=center] at (11.0,-1.0) {\texttt{My\$chool2024}\\Very hard to crack};
\end{tikzpicture}
\caption{Password strength meter --- aim for the green zone!}
\label{fig:password-meter}
\end{figure}

\begin{theorybox}[Creating a Strong Password]
A strong password should have:
\begin{itemize}[leftmargin=*, label={\textcolor{LabGreen}{\ding{51}}}, itemsep=2pt]
  \item At least \textbf{8 characters} (12+ is better)
  \item A mix of \textbf{uppercase} and \textbf{lowercase} letters
  \item At least \textbf{one number} (0--9)
  \item At least \textbf{one special character} (\!, @, \#, \$, \%, etc.)
  \item \textbf{No personal info} (not your name, birthday, or ``password'')
\end{itemize}
\end{theorybox}

\section{Common Threats}
\label{sec:threats}
\index{virus}\index{malware}\index{phishing}

\begin{itemize}[leftmargin=*, label={\textcolor{WarnRed}{\faExclamationTriangle}}, itemsep=4pt]
  \item \textbf{Virus}\index{virus} --- A program that spreads by infecting other files and can damage your system
  \item \textbf{Malware}\index{malware} --- Any malicious software (viruses, spyware, ransomware)
  \item \textbf{Phishing}\index{phishing} --- Fake emails or websites that trick you into sharing passwords
  \item \textbf{Spam}\index{spam} --- Unwanted junk messages, often containing scams or malware
\end{itemize}

\begin{warnbox}
Never click links in emails from unknown senders! Phishing emails often look like they're from real companies (banks, schools, social media). Always check the sender's email address carefully.
\end{warnbox}

\bytetip[warn]{Windows has a built-in security tool called \textbf{Windows Defender}. Make sure it's always turned on! It scans for viruses and blocks threats automatically. You can check it in Settings $\rightarrow$ Privacy \& Security $\rightarrow$ Windows Security.}

\begin{funfact}
The first computer virus was called \textbf{``Creeper''} and was created in 1971 --- just as an experiment! It displayed the message: ``I'm the creeper, catch me if you can!'' A program called ``Reaper'' was written to delete it. That makes Reaper the world's first antivirus!
\end{funfact}

\begin{reviewqs}
  \item What makes a strong password? List four characteristics.
  \item What is a computer virus?
  \item How is phishing different from a virus?
  \item Name the built-in security tool in Windows.
  \item Give three tips for staying safe online.
\end{reviewqs}

\begin{challenge}
Create \textbf{three strong passwords} using the rules you learned. Each must be at least 12 characters and include uppercase, lowercase, numbers, and symbols. Write them down (but don't use them for real accounts --- this is just practice!).
\end{challenge}
